% ejemplo.tex - apunte de referencia de Classmer: usa todo lo que trae classmer.cls.
%
% Lo único obligatorio es \documentclass{classmer} y los datos de la clase
% (materia, fecha, número, título). El resto es libre: paquetes, estructura,
% gráficos. Los entornos de la plantilla se usan cuando ayudan, no por obligación.
%
% Entornos (título opcional entre corchetes; se parten entre páginas):
%   definicion  teorema  ejemplo  intuicion  parcial  agregadoIA  duda
%   \nuevacaja{lema}{Lema}{teorema} crea otra caja con el look de una existente.
% Marcas: \ts{hh:mm:ss}   \pizarra{archivo}{hh:mm:ss}{epígrafe}
%         \nota{inaudible}  \nota{ilegible}  \nota{sic}
% Atajos: \R \Q \CC \Z \tq \ok \cross (si el apunte define los suyos, ganan esos).
% Colores: azul rojo verde gris violeta naranja, y fondoazul, fondoverde, ...
% Lo que no hay que tocar: \pagestyle y los entornos de arriba (son el look).
\documentclass{classmer}

% ---- Datos de la clase (salen de input/contexto.json) ----------------------
\materia{Análisis Matemático II}
\fechaclase{15 de octubre de 2026}
\numeroclase{7}
\tipoclase{Teórica}                    % opcional
\titulo{Derivada direccional y gradiente}

% ---- Libre: lo que el apunte necesite ---------------------------------------
\usepackage{pgfplots}                  % la clase ya le pone compat=1.18
\usetikzlibrary{decorations.pathreplacing}
\nuevacaja{lema}{Lema}{teorema}        % otra caja con el look de "teorema"

\begin{document}
% El bloque de título se imprime solo: no hace falta \maketitle.

\section{Repaso y motivación \ts{00:02:15}}

La clase arrancó retomando las derivadas parciales de la clase anterior
\ts{00:02:40}. El profe remarcó dos limitaciones:
\begin{itemize}
  \item las parciales sólo miran en las direcciones de los ejes;
  \item que existan todas las parciales no alcanza para que $f$ sea continua:
  \begin{itemize}
    \item el ejemplo de siempre es $f(x,y)=\frac{xy}{x^2+y^2}$ fuera del origen;
    \item en el origen las dos parciales valen $0$.
  \end{itemize}
\end{itemize}

\begin{intuicion}
Pararse en un punto de la superficie $z=f(x,y)$ y preguntar ``¿qué tan empinado
es si camino \emph{en esta dirección}?''. Las parciales responden sólo para dos
direcciones; la derivada direccional, para todas.
\end{intuicion}

\section{Derivada direccional \ts{00:34:10}}

\begin{definicion}[Derivada direccional]
Sean $U\subseteq\R^n$ abierto, $f\colon U\to\R$, $a\in U$ y $v\in\R^n$ con
$\|v\|=1$. La \textbf{derivada direccional} de $f$ en $a$ en la dirección $v$ es
\[
  \frac{\partial f}{\partial v}(a)=\lim_{t\to 0}\frac{f(a+tv)-f(a)}{t},
\]
si el límite existe.
\end{definicion}

\begin{definicion}
$v$ es una \textbf{dirección} si $\|v\|=1$. Con $v=e_i$ se recupera la derivada
parcial $\frac{\partial f}{\partial x_i}(a)$.
\end{definicion}

% La foto puede no estar: en ese caso sale un recuadro y la compilación sigue.
\pizarra{pizarra_07.jpg}{00:34:52}{La definición tal como quedó en la pizarra.}

\begin{ejemplo}[Una cuenta por definición]
$f(x,y)=x^2y$, $a=(1,2)$ y $v=\tfrac{1}{\sqrt2}(1,1)$:
\begin{align*}
  f(a+tv) &= \Bigl(1+\tfrac{t}{\sqrt2}\Bigr)^2\Bigl(2+\tfrac{t}{\sqrt2}\Bigr),\\
  \frac{\partial f}{\partial v}(a) &= \frac{d}{dt}\Big|_{t=0} f(a+tv)
     = \frac{4}{\sqrt2}+\frac{1}{\sqrt2} = \frac{5}{\sqrt2}.
\end{align*}
\end{ejemplo}

\begin{ejemplo}
Con $v=e_2$ queda $\frac{\partial f}{\partial y}(1,2)=1$ \ok. Con $v=(1,1)$
\cross: no es unitario, hay que normalizarlo primero.
\end{ejemplo}

\section{Gradiente \ts{00:41:00}}

\begin{teorema}[Derivada direccional y gradiente]
Si $f$ es diferenciable en $a$, para toda dirección $v$
\[
  \frac{\partial f}{\partial v}(a)=\langle\nabla f(a),v\rangle,
  \qquad
  \max_{\|v\|=1}\frac{\partial f}{\partial v}(a)=\|\nabla f(a)\|.
\]
\end{teorema}

\begin{proof}
Por la diferenciabilidad, $f(a+tv)=f(a)+t\langle\nabla f(a),v\rangle+o(t)$; se
divide por $t$ y se toma límite. El profe acotó algo sobre el resto
\nota{inaudible} y siguió con el máximo, que sale de Cauchy--Schwarz.
\end{proof}

\begin{teorema}
Si $f$ es diferenciable en $a$, entonces es continua en $a$ \ts{00:44:30}.
\end{teorema}

\begin{lema}[Una variable]
$g(t)=f(a+tv)$ cumple $g'(0)=\frac{\partial f}{\partial v}(a)$.
\end{lema}

El dibujo de la pizarra, pasado en limpio: las curvas de nivel de
$f(x,y)=x^2+y^2$ y el gradiente en un punto.
\begin{center}
\begin{tikzpicture}[>={Stealth[length=2mm]}, font=\small]
  \draw[->, gris] (-2.5,0) -- (2.8,0) node[below] {$x$};
  \draw[->, gris] (0,-2.3) -- (0,2.6) node[left] {$y$};
  \foreach \r in {0.7,1.3,1.9} \draw[azul] (0,0) circle (\r);
  \node[azul, below] at (-60:1.9) {$f=c$};
  \fill[rojo] (45:1.3) circle (1.6pt) node[below right] {$a$};
  \draw[->, thick, rojo] (45:1.3) -- (45:2.5) node[right] {$\nabla f(a)$};
  \draw[->, thick, verde] (45:1.3) -- ++(135:0.9)
    node[left, fill=white, inner sep=1pt] {$v\perp\nabla f(a)$};
  \draw[decorate, decoration={brace, amplitude=4pt, mirror}, gris]
    (0,0) -- (45:1.25) node[midway, below right=2pt, fill=white, inner sep=1pt]
    {$\|a\|$};
\end{tikzpicture}
\end{center}

\begin{duda}[Por qué el gradiente es perpendicular a la curva de nivel]
Nota de la clase \ts{00:41:03}: ``preguntar por qué el gradiente es
perpendicular''. El profe lo dio por visto en la pizarra \nota{ilegible}.
\end{duda}

\begin{duda}
No quedó claro si $\nabla f(a)\neq 0$ es hipótesis del teorema \ts{00:47:12}.
\end{duda}

\begin{agregadoIA}[Gráfico de $g(t)=f(a+tv)$ para el ejemplo]
No se mostró en clase. La pendiente en $t=0$ es la derivada direccional
$\frac{5}{\sqrt2}\approx 3{,}54$:
\begin{center}
\begin{tikzpicture}
  \begin{axis}[width=0.75\linewidth, height=4.6cm, axis lines=middle,
      xlabel={$t$}, ylabel={$g(t)$}, domain=-1.5:1.2, samples=60,
      xtick={-1.5,-1,-0.5,0.5,1}, ytick={2,4,6},
      grid=major, grid style={gris!20}, every axis plot/.append style={thick}]
    \addplot[azul] {(1+x/sqrt(2))^2*(2+x/sqrt(2))};
    \addplot[rojo, dashed, domain=-0.8:0.8] {2+5/sqrt(2)*x};
  \end{axis}
\end{tikzpicture}
\end{center}
\end{agregadoIA}

\begin{agregadoIA}
Si $v$ no es unitario, $\langle\nabla f(a),v\rangle$ queda multiplicado por
$\|v\|$: por eso la definición pide $\|v\|=1$.
\end{agregadoIA}

\section{Cierre y avisos \ts{00:52:00}}

% Origen visible: [clase] = marca ⭐ del usuario; [ia] = lo detectaste vos
% ("esto lo tomo", "ojo con esto", énfasis, repetición). Segundo opcional: título.
\begin{parcial}[ia][Primer parcial]
``Esto lo toman siempre'': calcular una derivada direccional por definición y
con el gradiente, y comparar los resultados \ts{00:52:30}.
\end{parcial}

\begin{parcial}[clase]
Entra hasta gradiente inclusive \ts{00:53:05}.
\end{parcial}

\begin{importante}[clase]
Si $f$ es diferenciable en $a$, entonces $D_vf(a)=\langle\nabla f(a),v\rangle$
\ts{00:41:10}.
\end{importante}

\begin{importante}[ia][La derivada direccional es un número]
El profe lo repitió tres veces: el gradiente es un vector; la derivada
direccional, un número \ts{00:47:40}.
\end{importante}

Para la próxima clase:
\begin{enumerate}
  \item hacer los ejercicios 1 a 6 de la práctica 4;
  \item repasar:
  \begin{enumerate}
    \item regla de la cadena;
    \item diferenciabilidad (``la definición de la clase pasada'' \nota{sic}).
  \end{enumerate}
\end{enumerate}

\begin{intuicion}[Para recordar]
El gradiente apunta hacia donde $f$ crece más rápido y mide cuánto crece.
\end{intuicion}

\end{document}
